% problem.json の入力（3企業×7労働者, f4・f5・w8 なし）を数式で書いたもの
\documentclass[a4paper,11pt]{ltjsarticle}
\usepackage{amsmath,amssymb}
\usepackage[margin=25mm]{geometry}
\pagestyle{empty}
\setlength{\parindent}{0pt}

\begin{document}

\section*{市場}

定員は
\[
  (q_{f_1}, q_{f_2}, q_{f_3}) = (2, 2, 2).
\]

企業の選好は
\begin{align*}
  f_1 &: w_2 \succ w_6 \succ w_7 \succ w_1 \succ w_5 \succ w_3, \\
  f_2 &: w_3 \succ w_1 \succ w_6 \succ w_2 \succ w_7 \succ w_4, \\
  f_3 &: w_7 \succ w_4 \succ w_2 \succ w_5 \succ w_1 \succ w_3 \succ w_6.
\end{align*}

労働者の選好は
\begin{align*}
  w_1 &: f_1 \succ f_3 \succ f_2, \\
  w_2 &: f_2 \succ f_1 \succ f_3, \\
  w_3 &: f_3 \succ f_2 \succ f_1, \\
  w_4 &: f_2 \succ f_3, \\
  w_5 &: f_1 \succ f_3, \\
  w_6 &: f_3 \succ f_2 \succ f_1, \\
  w_7 &: f_2 \succ f_1 \succ f_3.
\end{align*}

共通禁止ペア集合は
\[
  \mathcal{B} = \bigl\{ \{w_2, w_3\}, \{w_2, w_6\}, \{w_3, w_5\}, \{w_4, w_5\}, \{w_5, w_7\} \bigr\}.
\]

\end{document}
